\documentclass{article}
\usepackage{booktabs}
\usepackage{amsmath}
\usepackage[a4paper,margin=1in]{geometry}

\begin{document}

\begin{table}[!ht]
\caption{Maximum time-series length $n$ segmented within a $10$-second budget for the Gaussian and Poisson models, as a function of the number of true segments. Gaussian segments alternate means $0/10$ ($\sigma=1$); Poisson segments alternate rates $4/20$. BS = BinSeg (\texttt{fpop::multiBinSeg} for Gaussian, \texttt{changepoint::cpt.meanvar} for Poisson); OP and DUST are dust's own engines; PELT is \texttt{changepoint}'s reference implementation; FPOP/GFPOP are \texttt{fpop}/\texttt{gfpop}. All dust entries (OP, DUST) use the Highway (SIMD) backend. \label{tab:simulation_results}}
\begin{tabular*}{\columnwidth}{@{\extracolsep\fill}llccc@{\extracolsep\fill}}
\toprule
Model & algorithm & 10 segments & 100 segments & 1000 segments \\
\midrule
Gaussian & BS    & $175 \times 10^6$          & $37.0 \times 10^6$         & $\text{---}^\dagger$ \\
Gaussian & OP    & $197 \times 10^3$          & $197 \times 10^3$          & $197 \times 10^3$ \\
Gaussian & PELT  & $340 \times 10^3$          & $1.05 \times 10^6$         & $3.30 \times 10^6$ \\
Gaussian & FPOP  & $54.0 \times 10^6$         & $62.0 \times 10^6$         & $64.0 \times 10^6$ \\
Gaussian & DUST  & $\mathbf{81.0 \times 10^6}$ & $\mathbf{98.0 \times 10^6}$ & $\mathbf{115 \times 10^6}$ \\
\midrule
Poisson & BS    & $980 \times 10^{3\ddagger}$ & $980 \times 10^{3\ddagger}$ & $980 \times 10^{3\ddagger}$ \\
Poisson & OP    & $86.0 \times 10^3$         & $86.0 \times 10^3$         & $86.0 \times 10^3$ \\
Poisson & PELT  & $200 \times 10^3$          & $640 \times 10^3$          & $2.00 \times 10^6$ \\
Poisson & GFPOP & $2.70 \times 10^6$         & $3.40 \times 10^6$         & $4.40 \times 10^6$ \\
Poisson & DUST  & $\mathbf{39.0 \times 10^6}$ & $\mathbf{45.0 \times 10^6}$ & $\mathbf{55.0 \times 10^6}$ \\
\bottomrule
\end{tabular*}

\vspace{0.5em}
\footnotesize
$\dagger$ No stable value: at this scale, \texttt{multiBinSeg}'s running time is not a reproducible function of $n$ -- the identical $n$ ran in $1.7$s under one random draw and did not complete within $100$s under another, on freshly started processes (ruling out memory or cache effects). Elapsed time instead depends on the specific noise realization reaching \texttt{multiBinSeg}'s recursive splitting routine.

$\ddagger$ Crash-limited, not time-limited: \texttt{changepoint::cpt.meanvar(method="BinSeg")} overflows its C stack and crashes the R process above $n \approx 985\text{,}000$--$997\text{,}500$, independent of the number of segments; $980 \times 10^3$ is the largest value confirmed safe. The reported elapsed times at that $n$ are well under 10s (0.3s/2.1s/20.2s for 10/100/1000 segments respectively) and are not themselves the limiting factor.

\end{table}

\end{document}
